% !TEX program = xelatex
% AI-effects 프로젝트 --- 한국 직업(SOC)별 월간 업무용 AI 대화량 산출 과정.
% 표 본문: results/table/19_{korea,coverage,workpct,major,major_months,detailed}.tex
%          (code/04_analysis_kor_soc_conversations_summary.do가 생성)
% 데이터:  code/03_merge_sensortower_openai_messages.do, code/03_merge_sensortower_aei_kor_soc_conversations.do
% 컴파일: xelatex -> biber -> xelatex -> xelatex (kotex, biblatex/biber 필요)
\documentclass[11pt]{article}

\usepackage{fontspec}
\usepackage{kotex}
\usepackage[a4paper,margin=2.2cm]{geometry}
\usepackage{longtable}
\usepackage{booktabs}
\usepackage{array}
\usepackage{amsmath}
\usepackage{amssymb}
\usepackage{xcolor}
\usepackage{hyperref}
\usepackage{enumitem}
\usepackage{titlesec}
\usepackage[backend=biber,style=authoryear,maxcitenames=2,uniquelist=false,sorting=nyt]{biblatex}
\addbibresource{../../reference/references.bib}

\hypersetup{colorlinks=true,linkcolor=blue!50!black,citecolor=blue!50!black,urlcolor=blue!50!black}
\newcolumntype{L}[1]{>{\raggedright\arraybackslash}p{#1}}
\newcolumntype{R}[1]{>{\raggedleft\arraybackslash}p{#1}}
\setlength{\LTleft}{0pt}
\setlength{\LTright}{0pt}
\renewcommand{\arraystretch}{1.15}
\setlength{\tabcolsep}{4pt}

\title{한국 직업(SOC)별 월간 업무용 AI 대화량 산출\\[2mm]
\large --- Sensor Tower 기반 대화량 $\times$ AEI 한국 직업 비중(업무용 기준) ---}
\author{AI-effects 프로젝트}
\date{\today}

\begin{document}
\maketitle

\begin{abstract}
\noindent 본 문서는 report 17에서 만든 한국의 월간 업무용 AI 대화량을 미국 표준직업분류(SOC) 직업별로 나눈 과정을 기록한다. 직업 구성은 Anthropic Economic Index(AEI) release 2026-06-26\parencite{massenkoff-2026-cadences}의 한국 Claude.ai 직업별 대화 비중을 쓴다. 이 자료는 2026년 4월$\cdot$5월 두 달에 대해 SOC 대분류 22개(비중과 업무용 비중 등 모든 지표)와 세부 직업 389--424개(비중만)를 공개한다. 핵심은 \textbf{업무용 대화 비중의 적용}이다. 나누려는 대화량은 업무용 대화인데 AEI의 직업 비중은 업무$\cdot$개인$\cdot$학업 대화를 모두 포함하고, 직업마다 업무용 비중이 8\%에서 85\%까지 다르다. 그래서 대분류 비중에 그 직업의 업무용 비중을 곱해 다시 정규화한 ``업무용 대화 기준 비중''을 쓴다. 이 보정으로 컴퓨터$\cdot$수학의 비중은 22.8\%에서 27.7\%로 오르고, 교육$\cdot$도서관(12.9\% $\to$ 9.1\%), 판매, 지역사회$\cdot$사회서비스, 보건의료 전문$\cdot$기술은 내려간다. 대분류로 다시 계산한 한국 업무용 비중(2026년 5월 47.36\%)은 AEI가 공개한 한국 전체 값(47.24\%)과 0.12\%p 차이로 맞는다. 2026년 5월 한국 업무용 대화 3.46억 건(대화당 메시지 3건 기준) 가운데 컴퓨터$\cdot$수학 0.96억 건, 예술$\cdot$디자인$\cdot$미디어 0.66억 건, 사무$\cdot$행정 지원 0.38억 건이다.
\end{abstract}

\tableofcontents

% ============================================================
\section{목적과 개요}
% ============================================================

report 17은 Sensor Tower 이용자 수와 ChatGPT 1인당 메시지 수로 국가별 월간 업무용 대화량을 만들었다. 본 문서는 그중 한국의 값을 직업별로 나누어, 이 저장소의 노출도(exposure) 지표와 직업 단위로 연결할 수 있는 \textbf{직업$\times$월 단위의 업무용 AI 대화량}을 만든다. 구축 흐름은 다음과 같다.
\begin{enumerate}[leftmargin=*]
  \item \textbf{한국 업무용 대화량}(report 17): 이용자 수 $\times$ 1인당 메시지 $\times$ Signals 한국 업무 비중 $\div$ 대화당 메시지 수 $K$.
  \item \textbf{AEI 한국 직업 비중}: 대분류 비중(\texttt{pct})과 대분류별 업무용 비중(\texttt{use\_case\_work\_pct}), 세부 직업 비중.
  \item \textbf{업무용 대화 기준 대분류 비중}: \texttt{pct} $\times$ 업무용 비중을 합이 1이 되게 정규화.
  \item \textbf{세부 직업 비중}: 대분류 비중 $\times$ 대분류 안 세부 직업 비중.
  \item \textbf{직업별 대화량}: 한국 업무용 대화량 $\times$ 직업 비중.
\end{enumerate}

% ============================================================
\section{원자료}
% ============================================================

\begin{longtable}{L{3.0cm}L{5.4cm}L{7.4cm}}
\caption{원자료}\label{tab:raw}\\
\toprule
자료 & 파일 & 내용 \\
\midrule
\endfirsthead
\toprule
자료 & 파일 & 내용 \\
\midrule
\endhead
report 17 산출물 & \texttt{sensortower\_work\_messages\_by\_country\_monthly.dta} & South Korea 행의 \texttt{work\_messages}, \texttt{conv\_k1p5}--\texttt{conv\_k10}(2024-07--2026-05) \\
AEI release 2026-06-26 & \texttt{aei\_claude\_ai\_2026-06-26.csv} & \texttt{geo\_id} = KOR, \texttt{category\_name} = \texttt{soc\_occupation}. 2026-04, 2026-05 \\
AEI 국가별 정리본 & \texttt{aei\_claude\_ai\_country\_usage.dta} & 한국 전체 \texttt{use\_case\_work\_pct}(검증용, \texttt{01\_import\_aei\_country\_usage.do}) \\
\bottomrule
\end{longtable}

\subsection{한국 월간 업무용 대화량(report 17)}
report 17의 계산을 한국에 적용하면 다음과 같다. 8개 제품 True Audience 이용자 수\parencite{sensortower-2026-state-of-ai}에 이용자 1인당 월 메시지 53.75건을 곱해 월간 메시지를 만든다. 1인당 메시지는 ChatGPT 2025년 7월 전 세계 월 메시지\parencite{chatterji-2025-how-people-use-chatgpt}에 AEI 2025년 8월의 18개국 대화 비중 61.01\%\parencite{appel-2025-uneven-geographic-and-enterprise-ai}를 곱하고 18개국 ChatGPT 이용자로 나눈 값이다. 여기에 OpenAI Signals의 한국 ChatGPT 업무용 메시지 비중\parencite{openai-2026-signals-v2}을 곱해 업무용 메시지를, 대화당 메시지 수 $K \in \{1.5, 3, 5, 7, 10\}$로 나누어 업무용 대화량을 얻는다.

\begin{table}[h]
\centering
\caption{한국 월간 이용자$\cdot$메시지$\cdot$업무용 대화량}\label{tab:korea}
\begin{tabular}{lrr}
\toprule
항목 & 2026-04 & 2026-05 \\
\midrule
% 04_analysis_kor_soc_conversations_summary.do가 생성. 직접 고치지 말 것.
8개 제품 True Audience 이용자(백만 명) & 57.5 & 61.0 \\
이용자 1인당 월 메시지(건) & 53.75 & 53.75 \\
월간 메시지(백만 건) & 3089.1 & 3277.1 \\
업무용 비중(\%, Signals 한국 ChatGPT) & 33.6 & 31.7 \\
월간 업무용 메시지(백만 건) & 1037.9 & 1038.8 \\
\midrule
업무용 대화(백만 건), 대화당 메시지 1.5개 & 692.0 & 692.6 \\
업무용 대화(백만 건), 대화당 메시지 3개 & 346.0 & 346.3 \\
업무용 대화(백만 건), 대화당 메시지 5개 & 207.6 & 207.8 \\
업무용 대화(백만 건), 대화당 메시지 7개 & 148.3 & 148.4 \\
업무용 대화(백만 건), 대화당 메시지 10개 & 103.8 & 103.9 \\

\bottomrule
\end{tabular}
\end{table}

두 달의 업무용 대화량은 거의 같다. 이용자는 늘었지만(57.5 $\to$ 61.0백만 명) Signals 한국 업무 비중이 33.6\%에서 31.7\%로 낮아졌기 때문이다.

\subsection{AEI 한국 직업(SOC) 비중}
AEI release 2026-06-26은 Claude.ai(Free$\cdot$Pro$\cdot$Max, Cowork 포함) 대화를 국가$\times$월로 집계해 공개한다. \texttt{soc\_occupation} 범주는 대화가 해당하는 O*NET 과업을 미국 노동통계국의 SOC 직업으로 묶은 것이다. 국가 단위에서는 다음이 공개된다(데이터 문서의 Metric Availability).
\begin{itemize}
  \item \textbf{대분류}(\texttt{hierarchy\_level} = 1, Major Group, 2자리 코드): \texttt{pct}와 \texttt{use\_case\_work\_pct}를 포함한 모든 지표.
  \item \textbf{세부 직업}(\texttt{hierarchy\_level} = 0, O*NET-SOC 코드 예: 15-1252.00): \texttt{pct}만.
  \item \texttt{pct}는 그 국가 대화 중 해당 직업의 비중(\%), \texttt{use\_case\_work\_pct}는 그 직업 대화 중 업무용(use case = work) 비중(\%)이다. 값은 소수 둘째 자리로 반올림되어 있다.
  \item 표본 하한에 못 미친 셀은 공개되지 않으므로 공개 비중의 합은 100\%보다 작다(표 \ref{tab:coverage}).
\end{itemize}

\begin{table}[h]
\centering
\caption{AEI 한국 SOC 공개 범위}\label{tab:coverage}
\begin{tabular}{lrr}
\toprule
항목 & 2026-04 & 2026-05 \\
\midrule
% 04_analysis_kor_soc_conversations_summary.do가 생성. 직접 고치지 말 것.
대분류: 공개 수 & 22 & 22 \\
대분류: 공개 pct 합계(\%) & 97.37 & 98.94 \\
세부 직업: 공개 수 & 389 & 424 \\
세부 직업: 공개 pct 합계(\%) & 96.16 & 98.08 \\
세부 직업이 하나도 공개되지 않은 대분류 수 & 1 & 0 \\

\bottomrule
\end{tabular}
\end{table}

% ============================================================
\section{업무용 대화 비중의 적용}
% ============================================================

\subsection{왜 필요한가}
나누려는 대상은 한국의 \textbf{업무용} 대화량이다. 반면 AEI의 \texttt{pct}는 업무$\cdot$개인$\cdot$학업 대화를 모두 포함한 분포다. 직업마다 업무용 비중이 크게 다르므로(표 \ref{tab:major}의 업무용 비중 열: 지역사회$\cdot$사회서비스 8.06\%, 보건의료 전문$\cdot$기술 10.89\%, 음식 조리$\cdot$서빙 12.00\% 대 건설$\cdot$채굴 85.29\%, 건물$\cdot$부지 청소$\cdot$관리 71.54\%) \texttt{pct}를 그대로 쓰면 개인적 용도가 많은 직업(예: 건강 상담, 교육 질문)에 업무용 대화를 과대 배분하게 된다.

\subsection{식}
월 $t$, 대분류 $g$에 대해 AEI 대분류 비중 $p_{g,t}$(\texttt{pct})와 업무용 비중 $w_{g,t}$(\texttt{use\_case\_work\_pct})로 업무용 대화 기준 비중을 만든다.
\[
  s^{\text{work}}_{g,t} = \frac{p_{g,t}\, w_{g,t}}{\sum_{g'} p_{g',t}\, w_{g',t}}, \qquad
  s^{\text{all}}_{g,t} = \frac{p_{g,t}}{\sum_{g'} p_{g',t}} \ \text{(비교용)} .
\]
$p_{g,t} w_{g,t}/100$은 한국 전체 대화 중 ``직업 $g$이면서 업무용''인 대화의 비중이므로, 이를 정규화한 $s^{\text{work}}$는 한국 업무용 대화 안에서 직업 $g$이 차지하는 비중이다. 분모가 공개된 22개 대분류의 합이므로, 공개되지 않은 몫(1--3\%)은 공개된 직업에 비례 배분된다. 직업별 대화량은
\[
  \text{conv\_k}K_{g,t} = \text{conv\_k}K_{\text{KOR},t} \times s^{\text{work}}_{g,t}, \qquad K \in \{1.5, 3, 5, 7, 10\}
\]
이다. do 파일은 매월 직업별 대화량의 합이 한국 대화량과 같은지 검사한다.

\subsection{결과: 전체 기준 대 업무용 기준}
\begin{table}[htbp]
\centering
\small
\caption{SOC 대분류별 비중과 업무용 대화량(2026년 5월)}\label{tab:major}
\begin{tabular}{lrrrrrr}
\toprule
 & AEI & AEI & \multicolumn{2}{c}{비중(\%)} & 차이 & 업무용 대화 \\
\cmidrule(lr){4-5}
대분류 & pct(\%) & 업무용 비중(\%) & 전체 기준 & 업무용 기준 & (\%p) & ($K=3$, 백만 건) \\
\midrule
% 04_analysis_kor_soc_conversations_summary.do가 생성. 직접 고치지 말 것.
15 컴퓨터·수학 & 22.56 & 57.50 & 22.80 & 27.68 & 4.88 & 95.86 \\
27 예술·디자인·엔터테인먼트·스포츠·미디어 & 16.76 & 53.19 & 16.94 & 19.02 & 2.08 & 65.88 \\
43 사무·행정 지원 & 8.15 & 63.14 & 8.24 & 10.98 & 2.74 & 38.03 \\
25 교육·도서관 & 12.76 & 33.32 & 12.90 & 9.07 & -3.82 & 31.42 \\
11 경영 & 5.14 & 61.90 & 5.20 & 6.79 & 1.59 & 23.51 \\
13 사업·금융 운영 & 7.05 & 42.06 & 7.13 & 6.33 & -0.80 & 21.91 \\
17 건축·공학 & 3.72 & 62.83 & 3.76 & 4.99 & 1.23 & 17.27 \\
19 생명·물리·사회과학 & 5.34 & 40.54 & 5.40 & 4.62 & -0.78 & 16.00 \\
41 판매 & 7.00 & 26.95 & 7.07 & 4.03 & -3.05 & 13.94 \\
23 법률 & 0.98 & 60.38 & 0.99 & 1.26 & 0.27 & 4.37 \\
51 생산 & 0.89 & 64.05 & 0.90 & 1.22 & 0.32 & 4.21 \\
39 개인 돌봄·서비스 & 1.16 & 32.50 & 1.17 & 0.80 & -0.37 & 2.79 \\
49 설치·정비·수리 & 0.57 & 62.32 & 0.58 & 0.76 & 0.18 & 2.63 \\
29 보건의료 전문·기술 & 2.46 & 10.89 & 2.49 & 0.57 & -1.91 & 1.98 \\
21 지역사회·사회서비스 & 2.55 & 8.06 & 2.58 & 0.44 & -2.14 & 1.52 \\
53 운송·자재 이동 & 0.36 & 54.44 & 0.36 & 0.42 & 0.05 & 1.45 \\
33 보안 서비스 & 0.32 & 60.99 & 0.32 & 0.42 & 0.09 & 1.44 \\
31 보건의료 지원 & 0.59 & 14.38 & 0.60 & 0.18 & -0.42 & 0.63 \\
47 건설·채굴 & 0.08 & 85.29 & 0.08 & 0.15 & 0.06 & 0.50 \\
37 건물·부지 청소·관리 & 0.09 & 71.54 & 0.09 & 0.14 & 0.05 & 0.48 \\
35 음식 조리·서빙 & 0.38 & 12.00 & 0.38 & 0.10 & -0.29 & 0.34 \\
45 농림어업 & 0.03 & 58.33 & 0.03 & 0.04 & 0.01 & 0.13 \\
\midrule
합계 & 98.94 & 47.36 & 100.00 & 100.00 & 0.00 & 346.28 \\

\bottomrule
\end{tabular}
\par\smallskip\footnotesize 주: 업무용 기준 비중 순. 전체 기준 = $s^{\text{all}}$, 업무용 기준 = $s^{\text{work}}$, 차이 = 업무용 기준 $-$ 전체 기준. 합계 행의 업무용 비중은 함의된 한국 업무용 비중.
\end{table}

업무용 비중이 한국 평균(약 47\%)보다 높은 직업은 비중이 오르고 낮은 직업은 내려간다.
\begin{itemize}
  \item \textbf{오름}: 컴퓨터$\cdot$수학(22.80\% $\to$ 27.68\%, +4.88\%p), 사무$\cdot$행정 지원(+2.74\%p), 예술$\cdot$디자인$\cdot$미디어(+2.08\%p), 경영(+1.59\%p), 건축$\cdot$공학(+1.23\%p).
  \item \textbf{내림}: 교육$\cdot$도서관(12.90\% $\to$ 9.07\%, $-$3.82\%p), 판매($-$3.05\%p), 지역사회$\cdot$사회서비스($-$2.14\%p), 보건의료 전문$\cdot$기술($-$1.91\%p).
\end{itemize}
상위 3개 대분류(컴퓨터$\cdot$수학, 예술$\cdot$디자인$\cdot$미디어, 사무$\cdot$행정 지원)가 업무용 대화의 57.7\%를 차지한다.

\subsection{검증: 함의된 한국 업무용 비중}
대분류 값으로 한국 전체의 업무용 비중을 다시 계산하면 $\sum_g p_g w_g / \sum_g p_g$이다. 이 값은 AEI가 별도로 공개한 한국 전체(\texttt{overall}) 업무용 비중과 같아야 한다.

\begin{table}[h]
\centering
\caption{함의된 한국 업무용 비중 대 AEI 공개 값}\label{tab:workpct}
\begin{tabular}{lrr}
\toprule
항목 & 2026-04 & 2026-05 \\
\midrule
% 04_analysis_kor_soc_conversations_summary.do가 생성. 직접 고치지 말 것.
대분류로 함의된 업무용 비중(\%) = $\sum$ pct $\times$ 업무용 비중 / $\sum$ pct & 49.94 & 47.36 \\
AEI 공개 한국 전체 업무용 비중(\%) & 49.78 & 47.24 \\
차이(\%p) & 0.16 & 0.12 \\

\bottomrule
\end{tabular}
\end{table}

두 달 모두 차이가 0.2\%p 안이다. 남은 차이는 공개되지 않은 대분류 몫(1--3\%)과 반올림으로 설명된다. 따라서 대분류 \texttt{pct}와 \texttt{use\_case\_work\_pct}의 조합이 한국 전체 업무 대화를 일관되게 나눈다고 볼 수 있다.

\subsection{업무용 정의의 차이}
한국 업무용 대화의 \textbf{총량}은 Signals(ChatGPT, 메시지 단위 \texttt{work\_related} 이진 분류, 2026년 5월 31.7\%)로, \textbf{직업 구성}은 AEI(Claude.ai, 대화 단위 use case 3분류, 47.2\%)로 정한다. 두 정의는 분류기와 단위가 다르다(report 18). 본 문서는 AEI의 업무용 비중을 직업 간 \textbf{상대적} 차이로만 쓰고 총량에는 쓰지 않는다.

% ============================================================
\section{세부 직업 배분}
% ============================================================

세부 직업은 \texttt{pct}만 공개되므로 업무용 비중을 직접 적용할 수 없다. 세부 직업 $d$(대분류 $g$ 소속, O*NET-SOC 코드의 앞 2자리 = $g$)에 대해
\[
  s^{\text{in}}_{d,t} = \frac{p_{d,t}}{\sum_{d' \in g} p_{d',t}}, \qquad
  s^{\text{work}}_{d,t} = s^{\text{work}}_{g,t} \times s^{\text{in}}_{d,t}, \qquad
  \text{conv\_k}K_{d,t} = \text{conv\_k}K_{\text{KOR},t} \times s^{\text{work}}_{d,t}
\]
로 둔다. 이는 같은 대분류 안의 세부 직업은 업무용 비중이 같다는 가정(B3)이다. 대분류 안에서 공개되지 않은 세부 직업 몫은 공개된 세부 직업에 비례 배분된다. 2026년 4월 농림어업(45)은 공개된 세부 직업이 하나도 없어 \texttt{soc\_code} = ``45-XXXX.XX'', \texttt{detail\_missing} = 1인 한 행에 대분류 몫 전체를 둔다(대화량 합계 보존).

\begin{longtable}{lL{6.2cm}rrrrr}
\caption{세부 직업 상위 30개(2026년 5월, $K=3$ 업무용 대화 순)}\label{tab:detailed}\\
\toprule
 &  & AEI & 대분류 안 & 대분류 & 업무용 & 업무용 대화 \\
코드 & 직업 & pct(\%) & 비중(\%) & 비중(\%) & 비중(\%) & (백만 건) \\
\midrule
\endfirsthead
\toprule
코드 & 직업 & pct(\%) & 대분류 안(\%) & 대분류(\%) & 업무용(\%) & 백만 건 \\
\midrule
\endhead
% 04_analysis_kor_soc_conversations_summary.do가 생성. 직접 고치지 말 것.
15-1299.03 & Document Management Specialists & 5.49 & 24.30 & 27.68 & 6.73 & 23.30 \\
27-3041.00 & Editors & 3.27 & 19.56 & 19.02 & 3.72 & 12.88 \\
25-4022.00 & Librarians and Media Collections Specialists & 4.46 & 34.93 & 9.07 & 3.17 & 10.97 \\
27-3023.00 & News Analysts, Reporters, and Journalists & 2.60 & 15.55 & 19.02 & 2.96 & 10.24 \\
27-3043.00 & Writers and Authors & 2.56 & 15.31 & 19.02 & 2.91 & 10.09 \\
27-3042.00 & Technical Writers & 2.22 & 13.28 & 19.02 & 2.53 & 8.75 \\
15-1232.00 & Computer User Support Specialists & 1.90 & 8.41 & 27.68 & 2.33 & 8.06 \\
13-2071.00 & Credit Counselors & 2.45 & 34.85 & 6.33 & 2.21 & 7.64 \\
15-1211.00 & Computer Systems Analysts & 1.69 & 7.48 & 27.68 & 2.07 & 7.17 \\
27-3091.00 & Interpreters and Translators & 1.72 & 10.29 & 19.02 & 1.96 & 6.78 \\
27-3043.05 & Poets, Lyricists and Creative Writers & 1.60 & 9.57 & 19.02 & 1.82 & 6.30 \\
15-1255.00 & Web and Digital Interface Designers & 1.48 & 6.55 & 27.68 & 1.81 & 6.28 \\
15-1243.01 & Data Warehousing Specialists & 1.42 & 6.29 & 27.68 & 1.74 & 6.03 \\
15-1251.00 & Computer Programmers & 1.40 & 6.20 & 27.68 & 1.72 & 5.94 \\
25-4031.00 & Library Technicians & 2.30 & 18.01 & 9.07 & 1.63 & 5.66 \\
43-9022.00 & Word Processors and Typists & 1.15 & 14.27 & 10.98 & 1.57 & 5.43 \\
15-1253.00 & Software Quality Assurance Analysts and Testers & 1.26 & 5.58 & 27.68 & 1.54 & 5.35 \\
41-2021.00 & Counter and Rental Clerks & 1.83 & 26.26 & 4.03 & 1.06 & 3.66 \\
15-1299.08 & Computer Systems Engineers/Architects & 0.85 & 3.76 & 27.68 & 1.04 & 3.61 \\
15-1299.02 & Geographic Information Systems Technologists and Technicians & 0.85 & 3.76 & 27.68 & 1.04 & 3.61 \\
43-9061.00 & Office Clerks, General & 0.72 & 8.93 & 10.98 & 0.98 & 3.40 \\
17-2061.00 & Computer Hardware Engineers & 0.68 & 18.63 & 4.99 & 0.93 & 3.22 \\
15-2051.00 & Data Scientists & 0.75 & 3.32 & 27.68 & 0.92 & 3.18 \\
27-1024.00 & Graphic Designers & 0.74 & 4.43 & 19.02 & 0.84 & 2.92 \\
43-9081.00 & Proofreaders and Copy Markers & 0.60 & 7.44 & 10.98 & 0.82 & 2.83 \\
43-9111.00 & Statistical Assistants & 0.58 & 7.20 & 10.98 & 0.79 & 2.74 \\
23-1011.00 & Lawyers & 0.61 & 62.24 & 1.26 & 0.79 & 2.72 \\
43-9031.00 & Desktop Publishers & 0.57 & 7.07 & 10.98 & 0.78 & 2.69 \\
15-2021.00 & Mathematicians & 0.63 & 2.79 & 27.68 & 0.77 & 2.67 \\
43-6014.00 & Secretaries and Administrative Assistants, Except Legal, Medical, and Executive & 0.56 & 6.95 & 10.98 & 0.76 & 2.64 \\

\bottomrule
\end{longtable}

세부 직업 1위는 Document Management Specialists(15-1299.03, 업무용 비중 6.73\%, 0.23억 건)다. 상위 30개 중 편집자$\cdot$기자$\cdot$작가$\cdot$번역가 등 미디어 직업과 컴퓨터 직업이 대부분을 차지한다. 세부 직업 값은 대분류의 업무용 비중을 그대로 물려받으므로, 같은 대분류 안의 직업 간 순서는 \texttt{pct} 순서와 같다.

% ============================================================
\section{월별 비교}
% ============================================================

\begin{table}[htbp]
\centering
\small
\caption{대분류별 업무용 기준 비중과 대화량: 2026년 4월 대 5월}\label{tab:months}
\begin{tabular}{lrrrrr}
\toprule
 & \multicolumn{3}{c}{업무용 기준 비중(\%)} & \multicolumn{2}{c}{업무용 대화($K=3$, 백만 건)} \\
\cmidrule(lr){2-4}\cmidrule(lr){5-6}
대분류 & 2026-04 & 2026-05 & 차이(\%p) & 2026-04 & 2026-05 \\
\midrule
% 04_analysis_kor_soc_conversations_summary.do가 생성. 직접 고치지 말 것.
15 컴퓨터·수학 & 28.16 & 27.68 & -0.48 & 97.44 & 95.86 \\
27 예술·디자인·엔터테인먼트·스포츠·미디어 & 18.84 & 19.02 & 0.19 & 65.18 & 65.88 \\
43 사무·행정 지원 & 10.45 & 10.98 & 0.53 & 36.14 & 38.03 \\
25 교육·도서관 & 8.87 & 9.07 & 0.20 & 30.69 & 31.42 \\
11 경영 & 6.89 & 6.79 & -0.10 & 23.83 & 23.51 \\
13 사업·금융 운영 & 6.18 & 6.33 & 0.15 & 21.38 & 21.91 \\
17 건축·공학 & 5.21 & 4.99 & -0.22 & 18.03 & 17.27 \\
19 생명·물리·사회과학 & 4.68 & 4.62 & -0.06 & 16.20 & 16.00 \\
41 판매 & 4.40 & 4.03 & -0.37 & 15.22 & 13.94 \\
23 법률 & 1.11 & 1.26 & 0.15 & 3.85 & 4.37 \\
51 생산 & 1.28 & 1.22 & -0.06 & 4.41 & 4.21 \\
39 개인 돌봄·서비스 & 0.84 & 0.80 & -0.04 & 2.92 & 2.79 \\
49 설치·정비·수리 & 0.73 & 0.76 & 0.03 & 2.52 & 2.63 \\
29 보건의료 전문·기술 & 0.52 & 0.57 & 0.05 & 1.80 & 1.98 \\
21 지역사회·사회서비스 & 0.42 & 0.44 & 0.02 & 1.45 & 1.52 \\
53 운송·자재 이동 & 0.42 & 0.42 & 0.00 & 1.45 & 1.45 \\
33 보안 서비스 & 0.43 & 0.42 & -0.01 & 1.48 & 1.44 \\
31 보건의료 지원 & 0.14 & 0.18 & 0.04 & 0.48 & 0.63 \\
47 건설·채굴 & 0.11 & 0.15 & 0.04 & 0.38 & 0.50 \\
37 건물·부지 청소·관리 & 0.20 & 0.14 & -0.06 & 0.70 & 0.48 \\
35 음식 조리·서빙 & 0.08 & 0.10 & 0.02 & 0.28 & 0.34 \\
45 농림어업 & 0.04 & 0.04 & 0.00 & 0.13 & 0.13 \\

\bottomrule
\end{tabular}
\end{table}

두 달 사이 대분류 비중의 변화는 모두 $\pm$0.6\%p 안이다. 한 달 간격이라 직업 구성은 안정적이다. 컴퓨터$\cdot$수학($-$0.48\%p)과 판매($-$0.37\%p)가 줄고 사무$\cdot$행정 지원(+0.53\%p)이 늘었다.

% ============================================================
\section{산출 파일}
% ============================================================

\begin{longtable}{L{5.6cm}L{7.0cm}R{1.0cm}L{2.2cm}}
\caption{산출 파일(\texttt{data/proc/macro/scaling/sensortower\_state\_of\_ai\_2026/}, 각 .csv/.dta)}\label{tab:files}\\
\toprule
파일 & 주요 변수 & 행 수 & 키 \\
\midrule
\endfirsthead
\toprule
파일 & 주요 변수 & 행 수 & 키 \\
\midrule
\endhead
\texttt{sensortower\_aei\_kor\_soc\_major\_conversations} & \texttt{soc\_major}, \texttt{soc\_major\_title}, \texttt{pct}, \texttt{work\_pct}, \texttt{share\_all}, \texttt{share\_work}, \texttt{conv\_k*}, \texttt{work\_pct\_kor}, \texttt{kor\_conv\_k*} & 44 & month $\times$ soc\_major \\
\texttt{sensortower\_aei\_kor\_soc\_detailed\_conversations} & \texttt{soc\_code}, \texttt{soc\_title}, \texttt{soc\_major}, \texttt{pct}, \texttt{share\_in\_major}, \texttt{share\_work\_major}, \texttt{share\_work}, \texttt{conv\_k*}, \texttt{detail\_missing} & 814 & month $\times$ soc\_code \\
\bottomrule
\end{longtable}

기간은 AEI 직업 자료가 있는 2026-04, 2026-05 두 달이다. 요약표는 \texttt{results/table/19\_kor\_soc\_conversations.xlsx}(시트 \texttt{korea}, \texttt{major}, \texttt{detailed\_top30}, \texttt{coverage})에도 있다.

% ============================================================
\section{가정과 한계}
% ============================================================

\begin{enumerate}[leftmargin=*]
  \item \textbf{(B1) 제품 간 직업 구성 동일.} 8개 제품(ChatGPT, Gemini 등) 한국 업무용 대화의 직업 구성을 Claude.ai의 구성과 같다고 둔다. Claude는 코딩$\cdot$글쓰기 비중이 높은 제품이어서, 컴퓨터$\cdot$미디어 직업이 과대, 다른 직업이 과소 배분될 수 있다.
  \item \textbf{(B2) 미공개 셀 비례 배분.} 공개되지 않은 대분류 몫(1.1--2.6\%)과 세부 직업 몫은 공개된 셀과 같은 비율로 나뉜다고 둔다. 미공개 셀은 작은 직업이므로 작은 직업의 대화량이 다소 과소일 수 있다.
  \item \textbf{(B3) 대분류 안 업무용 비중 동일.} 세부 직업은 업무용 비중이 공개되지 않아 대분류 값을 쓴다.
  \item \textbf{SOC의 의미.} AEI의 직업은 사용자의 직업이 아니라 대화 내용이 해당하는 O*NET 과업의 직업이다. 예를 들어 Document Management Specialists의 대화량은 ``문서 관리 과업에 해당하는 업무용 대화''이지 그 직업 종사자의 대화가 아니다. 직업별 고용과 결합할 때 이 차이를 해석에 반영해야 한다.
  \item \textbf{업무용 정의의 혼합.} 총량은 Signals 정의, 구성은 AEI 정의를 쓴다(3.5절).
  \item \textbf{report 17 가정의 승계.} 한국 대화량 자체가 report 17의 (A1) 1인당 메시지 동일, (A2) 업무 비중 = ChatGPT, (A3) 18개국 비중 61.01\% 가정과 대화당 메시지 수 $K$ 격자에 의존한다. 직업별 값은 이 불확실성을 그대로 물려받는다.
  \item \textbf{기간.} 직업 자료가 두 달뿐이다. 다른 달로 넓히려면 직업 구성이 시기에 따라 같다는 가정이 추가로 필요하다.
  \item \textbf{반올림.} AEI 값은 소수 둘째 자리다. \texttt{pct}가 0.01--0.1\%인 세부 직업은 상대 오차가 크다.
\end{enumerate}

% ============================================================
\section{재현}
% ============================================================

Stata 17에서 다음 순서로 실행한다(모든 do 파일은 \texttt{\$root} 기준 경로).
\begin{enumerate}[leftmargin=*]
  \item \texttt{code/02\_clean\_sensortower\_share\_all\_products.do}, \texttt{code/03\_merge\_sensortower\_openai\_messages.do} --- report 17의 국가별 업무용 대화량.
  \item \texttt{code/01\_import\_aei\_country\_usage.do} --- AEI 국가별 정리본(검증용 한국 전체 업무용 비중).
  \item \texttt{code/03\_merge\_sensortower\_aei\_kor\_soc\_conversations.do} --- 한국 직업별 대화량 두 파일.
  \item \texttt{code/04\_analysis\_kor\_soc\_conversations\_summary.do} --- 본 문서의 표(\texttt{results/table/19\_*.tex}, \texttt{.xlsx}).
\end{enumerate}

\printbibliography[title={참고문헌}]

\end{document}
